% GENERATED_BY: run_oc_core_monograph_translation_rework_dispatch_v1.py
% AI_ASSISTED_TRANSLATION_DRAFT: TRUE
% THEOREM_REVIEW_STATUS: APPROVED
% THEOREM_REVIEW_MODE: FORMAL_AI_GATE__CODEX_CLI_CHATGPT
% THEOREM_REVIEW_REVIEWER_ID: CODEX_CLI_CHATGPT__AUTO_DISPATCH_20260406
% THEOREM_REVIEW_AUTHORITY_CLASS: FINAL_ADVERSARIAL_EXTERNAL_LLM_REVIEWER
% THEOREM_REVIEWED_AT: 2026-04-14T10:59:05Z
% SOURCE_LANGUAGE: EN
% TARGET_LANGUAGE: RU
% SOURCE_EN_REF: appendix/A_notation.tex
\section{Обозначения и символы}
\label{app:notation}

В этом приложении собраны основные символы и соглашения об обозначениях, используемые
во всей монографии OC Core~1.3. Оно не является исчерпывающим, но охватывает
структурные элементы, появляющиеся в большинстве глав.

\subsection{Структура континуума}

\begin{description}
  \item[\(K\)] Континуум (общий уровень).
  \item[\(K_x\)] Континуум на уровне \(x\in\{0,\dots,10\}\).
  \item[\(M_x\)] Пространство вложения для уровня \(K_x\).
  \item[\(\Omega(K)\)] Множество допустимых состояний \(K\).
  \item[\(\partial\Omega(K)\)] Граница допустимых состояний \(K\), определяемая
        насыщением порогов.
  \item[\(A(K)\)] Множество осей несовместимых различий:
        \(A(K) = \{A_1,\dots,A_n\}\).
  \item[\(P(t)\)] Вектор потенциалов (энергетических, химических, информационных,
        биологических, когнитивных, социальных и т.д.).
  \item[\(J(t)\)] Вектор потоков, преобразующих потенциалы и состояния.
  \item[\(\Theta(K)\)] Пороговый ландшафт \(K\), включающий пороги существования,
        устойчивости, критичности, размерности, смерти, выразительности и
        вложения.
  \item[\(C(K)\)] Семейство структурно устойчивых циклов \(K\).
  \item[\(k(K,t)\)] Мера континуумности (жизнеспособности) \(K\).
  \item[\(H(\Omega(K_x))\)] Историческое K-level обозначение булева фактора
        непустоты допустимой области; это не пороговый оператор
        \(H_{\mathrm{thr}}\).
  \item[\(\tau_{\mathrm{eff}}(K_x)\)] Скалярное эффективное время уровня.
        Если K-level глава локально пишет скалярное \(\tau(K_x)\), это
        читается как алиас \(\tau_{\mathrm{eff}}(K_x)\), а не как полное
        семейство временных параметров.
\end{description}

\subsection{Пороги и напряжение}

\begin{description}
  \item[\(\Theta_{\mathrm{exist}}\)] Пороги существования:
        минимальные условия для \(\Omega(K)\neq\emptyset\).
  \item[\(\Theta_{\mathrm{stab}}\)] Пороги устойчивости:
        условия ограниченной динамики.
  \item[\(\Theta_{\mathrm{crit}}\)] Критические пороги:
        поверхности качественного изменения или фазового перехода.
  \item[\(\Theta_{\mathrm{dim}}\)] Размерностные пороги:
        условия возникновения новых осей.
  \item[\(\Theta_{\mathrm{time}}\)] Порог временной устойчивости, используемый
        в K-level условиях возникновения времени; цикл даёт вклад в скалярное
        время уровня только при выполнении локального условия
        \(\Pi(C_j)>\Theta_{\mathrm{time}}\).
  \item[\(\Theta_{\mathrm{death}}\)] Пороги смерти:
        пределы, за которыми не остаётся допустимых состояний.
  \item[\(\Theta_{\mathrm{expr}}\)] Пороги выразительности:
        минимальная выразительная способность осей, необходимая для
        представления релевантных различий.
  \item[\(\Theta_{\mathrm{embed}}\)] Пороги вложения:
        ограничения, накладываемые пространствами вложения \(M_x\).
  \item[\(T(K,t)\)] Функционал структурного напряжения, зависящий от
        потенциалов, осей и градиентов.
\end{description}

\subsection{Операторы}

\begin{description}
  \item[\(E_{\mathrm{evol}}\)] Оператор структурной эволюции, который
        локально может писаться как \(E\), если глава явно связала этот
        символ:
        \(E_{\mathrm{evol}}: K(t) \to K(t+dt)\).
  \item[\(F\)] Оператор потоков:
        \(F: J(t)\mapsto J(t+dt)\).
  \item[\(G\)] Оператор потенциалов:
        \(G: P(t)\mapsto P(t+dt)\).
  \item[\(H_{\mathrm{thr}}\)] Пороговый оператор, который локально может
        писаться как \(H\), если глава явно связала этот символ:
        \(H_{\mathrm{thr}}: \Theta(t)\mapsto \Theta(t+dt)\).
  \item[\(Q\)] Оператор циклов:
        \(Q: C(t)\mapsto C(t+dt)\).
  \item[\(R\)] Оператор границы:
        \(R: \partial\Omega(t)\mapsto \partial\Omega(t+dt)\).
  \item[\(S_{\mathrm{axis}}\)] Оператор осевой структуры, который локально
        может писаться как \(S\), если глава явно связала этот символ:
        \(S_{\mathrm{axis}}: A(t)\mapsto A(t+dt)\).
  \item[\(U\)] Оператор континуумности:
        \(U: k(t)\mapsto k(t+dt)\).
  \item[\(\Psi_{x\to x+1}\)] Оператор рождения:
        переход от \(K_x\) к \(K_{x+1}\).
  \item[\(\Psi_{x\to y}\)] Обозначение не-соседнего маршрута рождения при
        \(y>x\). По умолчанию это не новый примитив, а сокращение допустимой
        композиции \(\Psi_{y-1\to y}\circ\cdots\circ\Psi_{x\to x+1}\), где все
        промежуточные пороговые и вложенные обязательства сохраняются.
  \item[\(E_{\mathrm{int}}\)] Оператор взаимодействия для нескольких континуумов.
\end{description}

\subsection{Уровни и пространства вложения}

\begin{description}
  \item[\(K_0\)] Структурный субстрат (без времени, без геометрии, без энергии).
  \item[\(K_1\)] Одномерные классические континуумы.
  \item[\(K_2\)] Физические континуумы (поля, фазы, перколяция, BKT,
        структуры, связанные с массой).
  \item[\(K_3\)] Химические континуумы (сети реакций, RAF-замыкание).
  \item[\(K_4\)] Протоклеточные континуумы (мембраны, градиенты, осмотические
        пороги и пороги кривизны).
  \item[\(K_5\)] Ранние нейронные и биоэлектрические континуумы (ионные каналы,
        возбудимость, спайки).
  \item[\(K_6\)] Когнитивные континуумы (репрезентации, связывание,
        предсказание).
  \item[\(K_7\)] Социальные континуумы (нормы, институты, доверие).
  \item[\(K_8\)] Цивилизационные континуумы (инфраструктуры, системные
        пороги).
  \item[\(K_9\)] Теоретические континуумы (теории, парадигмы, онтологии).
  \item[\(K_{10}\)] Мета--теоретические континуумы (моделирование моделей).
\end{description}

\subsection{Прочие обозначения}

\begin{description}
  \item[\(\dim(K)\)] Размерность континуума \(K\) (мощность или ранг
        множества его осей).
  \item[\(C_{\max}(K)\)] Максимальный структурно устойчивый комплекс циклов
        \(K\).
  \item[\(\mathrm{span}(A)\)] Линейная или структурная оболочка множества осей.
  \item[\(H_{\Omega}(K,t)\)] Индикатор существования в определении
        континуумности.
\end{description}

Это приложение с обозначениями намеренно компактно; более специализированные
символы, используемые только в отдельных работах-расширениях, определяются
локально в соответствующих текстах.
